\chapter{因果效应的定义}

作为人，你对因果推断的基本概念早已烂熟于心。仅仅因为活着，你就知道什么是因果效应，懂得关联与因果的区别，并且一生都在始终如一地运用这些知识。若非如此，你早已不在人世。没有基本的因果概念，你活不到读这一章的年纪，更谈不上学会阅读。蹒跚学步时，你看到那些跳进泳池的孩子后来够到了果酱罐，于是你也会一头扎进泳池；少年时，你看到那些冲下最危险坡道的人赢得了下一场滑雪赛，于是你也会从最陡的坡上冲下去；为人父母后，你观察到吃了药的孩子第二天没出现在公园里，于是你会拒绝给生病的孩子用抗生素。

既然你已经理解因果效应的定义、也明白关联与因果的区别，就不要指望从本章获得什么深刻的概念洞见。本章的目的，是引入一套数学记号，把你本就具备的因果直觉形式化。请务必把你自己的因果直觉与这里引入的记号对应起来。要精确地定义因果概念，这套记号是必需的，全书都将使用它。

\section{个体因果效应}

宙斯是一位等待心脏移植的病人。1 月 1 日，他接受了一颗新的心脏。五天后，他死了。设想我们能以某种方式知道——也许是靠神的启示——如果宙斯在 1 月 1 日没有接受心脏移植，他五天后仍会活着。掌握了这一信息，多数人会同意：移植导致了宙斯的死亡。心脏移植这一干预，对宙斯的五日存活产生了因果效应。

另一位病人赫拉，也在 1 月 1 日接受了心脏移植。五天后她活着。设想我们同样知道：如果赫拉在 1 月 1 日没有接受那颗心脏，她五天后也仍然活着。那么，移植对赫拉的五日存活没有因果效应。

这两个小故事说明了人是如何推理因果效应的：我们比较采取行动 $A$ 时的结局，与不采取行动 $A$ 时的结局（通常只在头脑中比较）。若两个结局不同，我们就说行动 $A$ 对结局有因果效应——或致病，或防病；否则，我们说行动 $A$ 对结局没有因果效应。

\mnote{“业”（karma）也是常用来指称导致结果的行动的词。}

流行病学家、统计学家、经济学家及其他社会科学家，把行动 $A$ 称为干预、暴露、政策或处理。

为使我们的因果直觉可以进行数学与统计分析，需要引入一些记号。考虑一个二值处理变量 $A$（1：受处理，0：未处理）与一个二值结局变量 $Y$（1：死亡，0：存活）。

\mnote{大写字母表示随机变量，小写字母表示随机变量的某个取值。}

本书把 $A$、$Y$ 这类对不同个体可取不同值的变量称为\emph{随机变量}。记 $Y^{a=1}$（读作“在治疗 $a=1$ 之下的 $Y$”）为在治疗取值 $a=1$ 时本可观测到的结局变量，$Y^{a=0}$（读作“在治疗 $a=0$ 之下的 $Y$”）为在治疗取值 $a=0$ 时本可观测到的结局变量。$Y^{a=1}$ 与 $Y^{a=0}$ 也是随机变量。宙斯有 $Y^{a=1}=1$、$Y^{a=0}=0$，因为他受治疗时会死、而未受治疗时本会活下来；赫拉则有 $Y^{a=1}=0$、$Y^{a=0}=0$，因为她受治疗时活了下来、未受治疗时也本会活下来。

\mnote{我们有时把“个体 $i$ 的结局为 $Y^{a}=1$”简写为 $Y_i^{a}=1$。严格说来，当 $i$ 指某个具体个体（如宙斯）时，$Y_i^{a}$ 不是随机变量，因为我们假定个体的反事实结局是确定性的（见技术注记 1.2）。}

于是可以给出个体因果效应的形式定义：若对某个个体有 $Y^{a=1}\neq Y^{a=0}$，则处理 $A$ 对该个体的结局 $Y$ 有因果效应。因此，处理对宙斯的结局有因果效应，因为 $Y^{a=1}=1\neq 0=Y^{a=0}$；但对赫拉的结局没有因果效应，因为 $Y^{a=1}=0=Y^{a=0}$。

变量 $Y^{a=1}$ 与 $Y^{a=0}$ 称为\emph{潜在结局}，或称\emph{反事实结局}。有些作者偏好“潜在结局”一词，以强调依所接受的处理不同，这两个结局中的任何一个都可能被观测到；另一些作者偏好“反事实结局”，以强调这些结局代表可能并未实际发生的情形（即与事实相反的情形）。

\begin{defbox}[个体 $i$ 的因果效应]
$Y_i^{a=1}-Y_i^{a=0}$
\end{defbox}

对每个个体而言，两个反事实结局中有一个——即与个体实际所接受的治疗取值相对应的那一个——是事实性的。例如，由于宙斯实际上受了治疗（$A=1$），他在治疗下的反事实结局 $Y^{a=1}=1$ 就等于他观测到的（实际的）结局 $Y=1$。也就是说，观测到的治疗为 $A=a$ 的个体，其观测结局 $Y$ 等于他的反事实结局 $Y^{a}$。这一等式可以简洁地写成 $Y=Y^{A}$，其中 $Y^{A}$ 表示把反事实 $Y^{a}$ 在个体实际治疗 $A$ 所对应的取值 $a$ 处取值。等式 $Y=Y^{A}$ 称为\emph{一致性}。

\begin{defbox}[一致性]
若 $A_i=a$，则 $Y_i^{a}=Y_i^{A}=Y_i$
\end{defbox}

个体因果效应被定义为反事实结局取值之间的对比，但对每个个体而言，这些结局中只有一个能被观测到——即与个体实际经历的治疗取值相对应的那个。其余反事实结局始终未被观测。由于数据缺失，个体效应不可识别，也就是说，无法把它表示为观测数据的函数（可能的例外见精要 2.1）。

\section{平均因果效应}

定义个体因果效应需要三条信息：所关心的结局、要比较的两个行动 $a=1$ 与 $a=0$，以及要比较其反事实结局 $Y^{a=0}$ 与 $Y^{a=1}$ 的那个个体。然而，由于识别个体因果效应一般不可能，我们转而关注一种聚合的因果效应：个体总体中的\emph{平均因果效应}。定义它同样需要三条信息：所关心的结局、要比较的行动 $a=1$ 与 $a=0$，以及一个定义明确的个体总体——要比较的是这个总体中诸个体的结局 $Y^{a=0}$ 与 $Y^{a=1}$。

以宙斯的大家族作为我们所关心的总体。表 1.1 给出了全部 20 位成员在接受治疗（$a=1$）与不接受治疗（$a=0$）两种情形下的反事实结局。看最后一列：这是每个个体若接受治疗（心脏移植）本可观测到的结局 $Y^{a=1}$。总体中有一半成员（20 人中的 10 人）若接受心脏移植本会死亡。也就是说，若总体中所有人都接受 $a=1$，本会发生该结局的个体比例为 $\Pr[Y^{a=1}=1]=10/20=0.5$。同样，由表 1.1 的另一列可以断定：总体中有一半成员（20 人中的 10 人）若未接受心脏移植本会死亡；即若总体中所有人都接受 $a=0$，本会发生该结局的个体比例为 $\Pr[Y^{a=0}=1]=10/20=0.5$。

我们通过清点死亡人数（10）并除以个体总数（20）算出了治疗之下的反事实风险为 0.5；这与在总体全部个体上对反事实结局取平均是同一回事。要看出二值结局下风险与平均的等价性，可用表 1.1 的数据计算 $Y^{a=1}$ 的平均值。

\begin{table}[htbp]
\centering
\caption{宙斯家族的潜在结局}
\label{tab:1.1}
\small
\begin{tabular}{lcc@{\hspace{2.5em}}lcc}
\toprule
& $Y^{a=0}$ & $Y^{a=1}$ & & $Y^{a=0}$ & $Y^{a=1}$\\
\midrule
Rheia      & 0 & 1 & Hephaestus & 0 & 1\\
Kronos     & 1 & 0 & Aphrodite  & 0 & 1\\
Demeter    & 0 & 0 & Polyphemus & 0 & 1\\
Hades      & 0 & 0 & Persephone & 1 & 1\\
Hestia     & 0 & 0 & Hermes     & 1 & 0\\
Poseidon   & 1 & 0 & Hebe       & 1 & 0\\
Hera       & 0 & 0 & Dionysus   & 1 & 0\\
Zeus       & 0 & 1 &            &   &  \\
Artemis    & 1 & 1 &            &   &  \\
Apollo     & 1 & 0 &            &   &  \\
Leto       & 0 & 1 &            &   &  \\
Ares       & 1 & 1 &            &   &  \\
Athena     & 1 & 1 &            &   &  \\
\bottomrule
\end{tabular}
\end{table}

现在可以给出总体中平均因果效应的形式定义：若在所关心的总体中，治疗之下的风险 $\Pr[Y^{a=1}=1]$ 与未治疗之下的风险 $\Pr[Y^{a=0}=1]$ 不相等，则称处理 $A$ 对结局 $Y$ 存在平均因果效应。

按此定义，处理 $A$ 在我们的总体中对结局 $Y$ 没有平均因果效应，因为治疗之下的死亡风险 $\Pr[Y^{a=1}=1]$ 与未治疗之下的死亡风险 $\Pr[Y^{a=0}=1]$ 都是 0.5。所有人都接受心脏移植、或所有人都不接受，结果并无分别：两种情形下都有一半人会死。当总体中的平均因果效应如本例这般为零时，我们说“无平均因果效应”这一零假设为真。由于风险等于平均值，又由于字母 E 通常用来表示总体平均或均值（也称“期望”），我们可以把总体中非零平均因果效应的定义改写为 $\E[Y^{a=1}]\neq\E[Y^{a=0}]$，从而使该定义同时适用于二值结局与非二值结局。

\begin{defbox}[总体中的平均因果效应]
$\E[Y^{a=1}]\neq\E[Y^{a=0}]$
\end{defbox}

“心脏移植 $A$ 的平均因果效应”是否存在，是由“接受心脏移植（$a=1$）”与“不接受心脏移植（$a=0$）”这两个行动之间的对比来定义的。当可能的行动多于两个（即处理不是二值的）时，特定的

\renewcommand{\FPname}{精要（Fine Point）}
\begin{finepoint}[干扰]
我们对反事实结局的定义隐含地假定：个体在治疗取值 $a$ 之下的反事实结局，不依赖于其他个体的治疗取值。例如，我们隐含地假定：无论赫拉是否也接受心脏移植，宙斯在接受心脏移植后都会死。也就是说，赫拉的治疗取值没有干扰宙斯的结局。反之，设想赫拉得到一颗新心脏这件事令宙斯心烦意乱，以致他挺不过自己的心脏移植手术——尽管若赫拉未接受移植，他本会活下来。在这种情形下，赫拉的治疗干扰了宙斯的结局。

在涉及传染性病原体或教育项目的研究中，个体之间的干扰很常见：在这类研究中，一个个体的结局受到他与其他总体成员之社会互动的影响。

存在干扰时，个体 $i$ 的反事实 $Y_i^{a}$ 定义不良，因为一个个体的结局依赖于其他个体的治疗取值。存在干扰时，“心脏移植对宙斯结局的因果效应”定义不良；人们需要转而谈论“在赫拉未得到新心脏时，心脏移植对宙斯结局的因果效应”，或“在赫拉得到新心脏时，心脏移植对宙斯结局的因果效应”。若其他亲友的治疗也干扰宙斯的结局，人们或许还需要谈论“在没有任何亲友得到新心脏时”“在只有赫拉得到新心脏时”等等情形下心脏移植对宙斯结局的因果效应，因为对每一特定的心脏分配方案，处理对宙斯结局的因果效应都可能不同。无干扰假设被 Cox (1958) 称为“单元之间无交互”，并被纳入 Rubin (1980) 所描述的“稳定单元治疗取值假设”（SUTVA）。关于干扰在因果效应定义中所扮演的角色，更详细的讨论见 Halloran and Struchiner (1995)、Sobel (2006)、Rosenbaum (2007) 与 Hudgens and Halloran (2009)。除另有说明外，本书通篇假定无干扰。
\end{finepoint}
\renewcommand{\FPname}{精要}

\renewcommand{\TPname}{技术注记（Technical Point）}
\begin{techpoint}[总体中的因果效应]
设 $\E[Y^{a}]$ 为总体中所有个体都接受治疗水平 $a$ 时的平均反事实结局。对离散结局，均值即期望值 $\E[Y^{a}]$ 定义为随机变量 $Y^{a}$ 一切可能取值 $y$ 上的加权和 $\sum_y y\,p_{Y^{a}}(y)$，其中 $p_{Y^{a}}(\cdot)$ 是 $Y^{a}$ 的概率质量函数，即 $p_{Y^{a}}(y)=\Pr[Y^{a}=y]$。对二值结局，$\E[Y^{a}]=\Pr[Y^{a}=1]$。对连续结局，期望值 $\E[Y^{a}]$ 定义为随机变量 $Y^{a}$ 一切可能取值 $y$ 上的积分 $\int y f_{Y^{a}}(y)\,dy$，其中 $f_{Y^{a}}(\cdot)$ 是 $Y^{a}$ 的概率密度函数。同时适用于离散与连续结局的一种常见表示是 $\E[Y^{a}]=\int y\,dF_{Y^{a}}(y)$，其中 $F_{Y^{a}}(\cdot)$ 是随机变量 $Y^{a}$ 的累积分布函数（cdf）。若对任意两个取值 $a$ 与 $a'$ 都有 $\E[Y^{a}]\neq\E[Y^{a'}]$，则称总体中存在非零平均因果效应。

由反事实结局之均值对比所定义的平均因果效应，是最常用的总体因果效应。然而，总体因果效应也可以定义为反事实结局之泛函（包括中位数、方差、风险率或 cdf）的对比。一般地，总体因果效应可以定义为不同行动或治疗取值下反事实结局边缘分布之任一泛函的对比。例如，总体对风险的因果效应定义为 $\Var(Y^{a=1})-\Var(Y^{a=0})$。
\end{techpoint}
\renewcommand{\TPname}{技术注记}

效应对比必须明确指定。例如，“阿司匹林的因果效应”若不说明所关心的对比是什么，就是一句无意义的话：我们必须指明，所关心的对比是“在世期间每日口服（必要时经鼻胃管）150 毫克阿司匹林、持续 5 年”与“不服用阿司匹林”之间的对比。即便在其他干预（例如“在世期间每日经皮肤吸收 500 毫克阿司匹林、持续 5 年”）下的反事实结局定义不良或根本不存在，这一因果效应仍然定义良好。

没有平均因果效应，并不意味着没有个体效应。表~1.1 表明，处理对总体中 12 名成员（包括宙斯）有因果效应，因为这 12 人的反事实结局 $Y^{a=1}$ 与 $Y^{a=0}$ 取值不同。这 12 人中，6 人因治疗而受害（$Y^{a=1}-Y^{a=0}=1$），6 人因治疗而受益（$Y^{a=1}-Y^{a=0}=-1$）。总体中的平均因果效应为

\begin{defbox}[总体中的平均因果效应]

\[\E[Y^{a=1}]-\E[Y^{a=0}]=\frac{6\times 1+6\times(-1)}{20}=0\]

\end{defbox}

这一相等绝非巧合：平均因果效应 $\E[Y^{a=1}]-\E[Y^{a=0}]$ 恒等于个体因果效应 $Y^{a=1}-Y^{a=0}$ 的平均值 $\E[Y^{a=1}-Y^{a=0}]$，因为平均值之差等于差值的平均值。当总体中每个个体都没有因果效应，即对所有个体都有 $Y^{a=1}=Y^{a=0}$ 时，我们说\emph{严格因果零假设}为真。严格因果零假设蕴含“无平均因果效应”这一零假设。

如后续各章所述，平均因果效应有时可以从数据中识别，即便个体因果效应不能。此后我们把“平均因果效应”简称为“因果效应”，把“无平均因果效应的零假设”简称为因果零假设。下面我们描述度量因果效应大小的不同方法。

\section{因果效应的度量}

我们已经看到，处理“心脏移植”$A$ 对结局“死亡”$Y$ 在宙斯家族这 20 人构成的总体中没有因果效应。因果零假设之所以成立，是因为治疗之下的反事实风险 $\Pr[Y^{a=1}=1]$ 与未治疗之下的反事实风险 $\Pr[Y^{a=0}=1]$ 都等于 0.5。表示因果零效应有几种等价方式。例如，我们可以说风险 $\Pr[Y^{a=1}=1]$ 减去风险 $\Pr[Y^{a=0}=1]$ 等于零（$0.5-0.5=0$），或者说风险 $\Pr[Y^{a=1}=1]$ 除以风险 $\Pr[Y^{a=0}=1]$ 等于一（$0.5/0.5=1$）。也就是说，因果零效应可以表示为
\begin{alignat*}{2}
\text{(i)}\quad && \Pr[Y^{a=1}=1]-\Pr[Y^{a=0}=1] &= 0\\
\text{(ii)}\quad && \frac{\Pr[Y^{a=1}=1]}{\Pr[Y^{a=0}=1]} &= 1\\
\text{(iii)}\quad && \frac{\Pr[Y^{a=1}=1]/\Pr[Y^{a=1}=0]}{\Pr[Y^{a=0}=1]/\Pr[Y^{a=0}=0]} &= 1
\end{alignat*}
其中等式 (i)、(ii)、(iii) 的左端分别是因果风险差、因果风险比与因果优势比。

现在设想另一种处理 $A$——吸烟，它对另一种结局 $Y$——肺癌——在我们的总体中有因果效应。此时因果零假设不成立：$\Pr[Y^{a=1}=1]$ 与 $\Pr[Y^{a=0}=1]$ 不相等。在这种情形下，因果风险差、风险比与优势比分别不为 0、1 与 1。相反，这些因果参数在不同的尺度上量化同一因果效应的强度。由于因果风险差、风险比与优势比（以及其他概括量）都在度量因果效应，我们称它们为\emph{效应度量}。另见精要 1.2。

\mnote{总体中的因果风险差是个体因果效应 $Y^{a=1}-Y^{a=0}$ 在差值尺度上的平均，即它是个体因果效应的平均值度量。相反，总体中的因果风险比并不是个体因果效应 $Y^{a=1}/Y^{a=0}$ 在比值尺度上的平均：它度量的是总体中的因果效应，但不是任何个体因果效应的平均。}

每种效应度量各有其用。例如，设想一个很大的总体：若接受治疗，百万人中有 3 人会发病；若未接受治疗，百万人中有 1 人会发病。因果风险比为 3，因果风险差为 0.000002。因果风险比（乘法尺度）用来计算相对于不治疗，治疗使发病风险增加多少倍；因果风险差（加法尺度）用来计算可归因于治疗的绝对病例数。究竟采用乘法尺度还是加法尺度，取决于推断的目的。

\begin{finepoint}[需治疗人数]
考虑一个有一亿病人的总体：若接受治疗（$a=1$），其中两千万人会在五年内死亡；若未接受治疗（$a=0$），其中三千万人会在五年内死亡。这一信息可以用几种等价的方式概括：
\begin{itemize}[leftmargin=1.4em,itemsep=0.1em,topsep=0.2em]
\item 因果风险差为 $\Pr[Y^{a=1}=1]-\Pr[Y^{a=0}=1]=0.2-0.3=-0.1$；
\item 若治疗这一亿病人，将比不治疗少死一千万人；
\item 需要治疗一亿病人，才能挽救一千万条生命；
\item 平均而言，需要治疗 10 人才能挽救 1 条生命。
\end{itemize}
我们把为减少一例 $Y=1$ 而平均需要接受治疗 $a=1$ 的人数，称为\emph{需治疗人数}（number needed to treat，NNT）。在上例中，NNT 等于 10。对于减少平均病例数的治疗（即因果风险差为负），NNT 等于因果风险差绝对值的倒数：
\[
\mathrm{NNT}=\frac{-1}{\Pr[Y^{a=1}=1]-\Pr[Y^{a=0}=1]}.
\]
对于增加平均病例数的治疗（即因果风险差为正），可以对称地定义\emph{需伤害人数}。NNT 由 Laupacis、Sackett 与 Roberts (1988) 引入。与因果风险差一样，NNT 只适用于其所依据的总体与时间区间。关于 NNT 作为效应度量的利弊，见 Grieve (2003)。
\end{finepoint}

\section{随机变异}

读到这里你或许会抱怨：我们计算效应度量的做法有点不近情理。我们不仅忽略了“不朽的宙斯不会死”这一众所周知的事实，更要紧的是——表 1.1 中的总体只有 20 个人。而我们关心的总体通常要大得多。

在我们这个小小的总体中，我们收集了全部个体的信息。在实践中，研究者只收集总体中一个样本的信息。即便所有研究个体的反事实结局都已知，依据样本工作也无法得到总体中在治疗取值 $a$ 之下发生该结局的个体的确切比例，即未治疗之下的死亡概率 $\Pr[Y^{a=0}=1]$ 无法直接计算。人们只能估计这一概率。

考虑表 1.1 中的个体。此前我们把它们视为一个 20 人的总体。现在设想把它们视为从一个远大于此、近乎无限的超总体（例如所有不朽者）中随机抽取的样本。我们用 $\widehat{\Pr}[Y^{a=0}=1]=10/20=0.50$ 表示样本中若未暴露本会死亡的个体比例。样本比例 $\widehat{\Pr}[Y^{a=0}=1]$ 不必恰好等于整个超总体若未暴露本会死亡的个体比例 $\Pr[Y^{a=0}=1]$。例如，设总体中 $\Pr[Y^{a=0}=1]=0.57$，但由于抽样变异带来的随机误差，我们这个小样本中 $\widehat{\Pr}[Y^{a=0}=1]=0.5$。

\mnote{随机误差的第一个来源：抽样变异}

我们用样本比例 $\widehat{\Pr}[Y^{a}=1]$ 来估计治疗取值 $a$ 之下的超总体概率 $\Pr[Y^{a}=1]$。$\Pr$ 上方的“帽号”表明样本比例 $\widehat{\Pr}[Y^{a}=1]$ 是相应总体量 $\Pr[Y^{a}=1]$ 的估计量。我们说 $\widehat{\Pr}[Y^{a}=1]$ 是 $\Pr[Y^{a}=1]$ 的\emph{一致估计量}，因为样本中个体数越多，$\widehat{\Pr}[Y^{a}=1]$ 与 $\Pr[Y^{a}=1]$ 之间的差异就越小。这是因为抽样变异所致的误差是随机的，因而服从大数定律。

\mnote{若当样本量趋于无穷时，差值 $\hat\theta-\theta$ 以趋于 1 的概率趋于零，则称 $\theta$ 的估计量 $\hat\theta$ 为一致估计量。}

由于超总体概率 $\Pr[Y^{a}=1]$ 无法计算，只能由样本比例 $\widehat{\Pr}[Y^{a}=1]$ 一致地估计，我们无法确切断定因果效应存在或不存在。相反，必须借助统计程序来评估关于因果零假设 $\Pr[Y^{a=1}=1]=\Pr[Y^{a=0}=1]$ 的经验证据（细节见第 10 章）。

\mnote{注意：把“一致性”用于估计量时，含义与用于反事实时不同。}

到目前为止，我们只把抽样变异当作随机误差的一个来源。但还可能有另一个随机变异来源：个体的反事实结局取值或许并非事先固定。我们把反事实结局 $Y^{a}$ 定义为个体若接受治疗取值 $a$ 时的结局。例如在第一个小故事中，宙斯若受治疗会死、若未受治疗本会活下来。按此定义，反事实结局的取值对每个个体都是固定的、确定性的，即宙斯有 $Y^{a=1}=1$、$Y^{a=0}=0$。换言之，宙斯若受治疗有 100\% 的概率死亡，若未受治疗有 0\% 的概率死亡。

\mnote{随机误差的第二个来源：非确定性反事实}

然而，我们可以设想另一种情形：宙斯若受治疗有 90\% 的概率死亡，若未受治疗有 10\% 的概率死亡。在这种情形下，反事实结局是随机的、非确定性的，因为宙斯在治疗下（0.9）与未治疗下（0.1）的死亡概率既不为零也不为一。表 1.1 所示的 $Y^{a=1}$ 与 $Y^{a=0}$ 的取值，将是以这些概率“抛掷生死硬币”的可能实现。此外，人们会预期这些概率因个体而异，因为并非所有个体都同样易感于该结局。与经典力学不同，量子力学认为结局本质上是非确定性的：即若宙斯死亡的量子力学概率为 90\%，则该理论认为，无论我们收集多少关于宙斯的数据，宙斯若受治疗是否真会发病的不确定性都是不可约的。

因此，在因果推断中，随机误差来自抽样变异、非确定性反事实，或二者兼有。然而，出于教学上的考虑，在第 10 章之前我们将基本忽略随机误差。具体而言，我们将假定反事实结局是确定性的，并假定我们已记录了一个非常大的（也许是假设的）超总体中每个个体的数据。这相当于把我们这个 20 人的总体视为一个 200 亿人的总体，其中有 10 亿人与宙斯相同，10 亿人与赫拉相同，如此等等。因此，在第 10 章之前，我们将以奥林匹斯山式的确定性进行计算。

随后在第 10 章我们将说明：无论个体层面上的世界是随机的（量子）还是确定性的（经典），超总体中因果效应的统计估计与置信区间都是相同的；相反，实际研究样本中平均因果效应的置信区间则会因反事实是确定性还是随机性而不同。所幸，超总体效应在多数情形下才是真正有实质意义的因果效应。

\begin{techpoint}[非确定性反事实]
对非确定性反事实结局，治疗取值 $a$ 之下的平均结局 $\E[Y^{a}]$ 等于随机变量 $Y^{a}$ 一切可能取值 $y$ 上的加权和 $\sum_y y\,p_{Y^{a}}(y)$，其中概率质量函数 $p_{Y^{a}}(\cdot)=\E[Q_{Y^{a}}(\cdot)]$，而 $Q_{Y^{a}}(y)$ 是在治疗水平 $a$ 之下发生结局 $Y=y$ 的随机概率。在正文所述的例子中，宙斯的 $Q_{Y^{a=1}}(1)=0.9$。（对连续结局，加权和改为积分。）

更一般地，非确定性的反事实结局定义并不给每个个体指定随机变量 $Y^{a}$ 的某个具体取值，而是指定 $Y^{a}$ 的一个个体特定的统计分布 $\Theta_{Y^{a}}(\cdot)$。非确定性的因果效应定义是确定性定义的推广，其中 $\Theta_{Y^{a}}(\cdot)$ 现在是一个可在 0 与 1 之间取值的随机 cdf。总体中的平均反事实结局 $\E[Y^{a}]=\E\bigl[\E\{Y^{a}\mid\Theta_{Y^{a}}(\cdot)\}\bigr]$。因此，
\[
\E[Y^{a}]=\E\!\left[\int y\,d\Theta_{Y^{a}}(y)\right]=\int y\,d\,\E[\Theta_{Y^{a}}(y)]=\int y\,dF_{Y^{a}}(y),
\]
其中 $F_{Y^{a}}(\cdot)=\E[\Theta_{Y^{a}_i}(\cdot)]$。

若反事实结局是二值且非确定性的，则总体中的因果风险比 $\E[Q_{Y^{a=1}}(1)]/\E[Q_{Y^{a=0}}(1)]$ 等于比值尺度上个体因果效应 $Q_{Y^{a=1}}(1)/Q_{Y^{a=0}}(1)$ 的加权平均 $\E[W\{Q_{Y^{a=1}}(1)/Q_{Y^{a=0}}(1)\}]$，权重为 $W=Q_{Y^{a=0}}(1)/\E[Q_{Y^{a=0}}(1)]$，前提是对总体中任何人都有 $Q_{Y^{a=0}}(1)\neq 0$（即确定性）。
\end{techpoint}

\section{因果与关联}

\mnote{关联与因果之别，正是全书其余部分要展开的主题。}

显然，实际研究可得的数据与表 1.1 所示的大不相同。例如，我们通常不能既知道宙斯受治疗时的结局 $Y^{a=1}$、又知道宙斯未受治疗时的结局 $Y^{a=0}$。在真实世界中，我们只能观测到其中一个结局，因为宙斯或者受治疗、或者未受治疗。我们把观测到的结局记作 $Y$。因此，对每个个体而言，我们知道其观测到的治疗水平 $A$ 与结局 $Y$，如表 1.2 所示。

\begin{table}[htbp]
\centering
\caption{观测数据（按治疗组汇总；原书按个体逐行列出 20 人）}
\label{tab:1.2}
\small
\begin{tabular}{lccc}
\toprule
治疗组 & 人数 & 死亡数（$Y=1$） & 死亡风险 \\
\midrule
$A=1$（受治疗） & 13 & 7 & $7/13$ \\
$A=0$（未受治疗） & 7 & 3 & $3/7$ \\
\bottomrule
\end{tabular}
\end{table}

表 1.2 的数据可用来计算总体中“恰好接受治疗取值 $a$”的那部分个体中发生结局 $Y$ 的比例。例如，在表 1.2 中，13 位受治疗（$A=1$）的个体里有 7 人死亡（$Y=1$）。因此，受治疗者的死亡风险 $\Pr[Y=1\mid A=1]$ 为 $7/13$。更一般地，条件概率 $\Pr[Y=1\mid A=a]$ 定义为：在所关心的总体中，恰好接受治疗取值 $a$ 的那部分个体里，发生结局 $Y$ 的个体所占的比例。

当受治疗者中发生结局的比例 $\Pr[Y=1\mid A=1]$ 等于未受治疗者中发生结局的比例 $\Pr[Y=1\mid A=0]$ 时，我们说处理 $A$ 与结局 $Y$ \emph{独立}、$A$ 与 $Y$ 无关联，或 $A$ 不能预测 $Y$。独立性记作 $Y\indep A$——等价地，$A\indep Y$——读作“$Y$ 与 $A$ 独立”。独立性的几种等价定义是
\begin{alignat*}{2}
\text{(i)}\quad && \Pr[Y=1\mid A=1]-\Pr[Y=1\mid A=0] &= 0\\
\text{(ii)}\quad && \frac{\Pr[Y=1\mid A=1]}{\Pr[Y=1\mid A=0]} &= 1\\
\text{(iii)}\quad && \frac{\Pr[Y=1\mid A=1]/\Pr[Y=0\mid A=1]}{\Pr[Y=1\mid A=0]/\Pr[Y=0\mid A=0]} &= 1
\end{alignat*}
其中 (i)、(ii)、(iii) 的左端分别是关联风险差、关联风险比与关联优势比。

\mnote{Dawid (1979) 引入符号 $\indep$ 表示独立。}

我们说处理 $A$ 与结局 $Y$ 有关联、或相互依赖，是指 $\Pr[Y=1\mid A=1]\neq\Pr[Y=1\mid A=0]$。在我们的总体中，处理与结局确有关联，因为受治疗者的死亡风险 $\Pr[Y=1\mid A=1]=7/13=0.54$ 大于未受治疗者的死亡风险 $\Pr[Y=1\mid A=0]=3/7=0.43$。关联风险差、风险比与优势比（以及其他度量）量化了关联存在时的强度。它们在不同的尺度上度量关联，我们称其为\emph{关联度量}。这些度量同样受随机变异的影响。不过，在第 10 章之前，我们将假定表 1.2 中的总体极大，从而忽略统计问题。

\mnote{对连续结局 $Y$，我们把处理与结局之间的\emph{均值独立}定义为：$\E[Y\mid A=1]=\E[Y\mid A=0]$。对二值结局，独立与均值独立是同一个概念。}

对二值结局，风险等于平均值，因此可以把总体中关联的定义改写为 $\E[Y\mid A=1]\neq\E[Y\mid A=0]$。对连续结局 $Y$，我们也把关联定义为 $\E[Y\mid A=1]\neq\E[Y\mid A=0]$。对二值 $A$、$Y$，$A$ 与 $Y$ 无关联，当且仅当二者不统计相关。

在我们这个 20 人的总体中，我们发现了 (i) 比较全部 20 人都受治疗时的死亡风险与全部 20 人都未受治疗时的死亡风险，没有因果效应；(ii) 比较恰好受治疗的 13 人的死亡风险与恰好未受治疗的 7 人的死亡风险，存在关联。图 1.1 描绘了因果与关联之别。总体（用一个菱形表示）被划分为白色区域（受治疗者）与较小的灰色区域（未受治疗者）。

\begin{figure}[htbp]
\centering
\begin{tikzpicture}[causal,every node/.style={transform shape},scale=0.9]
  % 总体：菱形，左白（受治疗）右灰（未受治疗）
  \begin{scope}
    \clip (0,0.95)--(-1.25,0)--(0,-0.95)--(1.25,0)--cycle;
    \fill[black!28] (0,-1.1) rectangle (1.3,1.1);
    \fill[white] (-1.3,-1.1) rectangle (0,1.1);
  \end{scope}
  \draw[thick] (0,0.95)--(-1.25,0)--(0,-0.95)--(1.25,0)--cycle;
  \draw[thick] (0,0.95)--(0,-0.95);
  \node[font=\small,above=3mm] at (0,0.95) {所关心的总体};
  \node[font=\small,anchor=east] at (-1.4,0) {受治疗};
  \node[font=\small,anchor=west] at (1.4,0) {未受治疗};
  % 箭头
  \draw[->,semithick] (-0.45,-1.0) -- (-1.3,-1.72);
  \draw[->,semithick] (0.45,-1.0) -- (1.3,-1.72);
  \node[font=\small] at (-1.45,-2.0) {因果};
  \node[font=\small] at (1.45,-2.0) {关联};
  % 因果：两个完整菱形
  \begin{scope}[shift={(-1.45,-3.05)}]
    \fill[white] (0,0.62)--(-0.8,0)--(0,-0.62)--(0.8,0)--cycle;
    \draw[thick] (0,0.62)--(-0.8,0)--(0,-0.62)--(0.8,0)--cycle;
  \end{scope}
  \begin{scope}[shift={(1.45,-3.05)}]
    \fill[black!28] (0,0.62)--(-0.8,0)--(0,-0.62)--(0.8,0)--cycle;
    \draw[thick] (0,0.62)--(-0.8,0)--(0,-0.62)--(0.8,0)--cycle;
  \end{scope}
  \node[font=\footnotesize] at (0,-3.05) {vs.};
  \node[font=\footnotesize] at (-1.45,-3.95) {$\E[Y^{a=1}]$};
  \node[font=\footnotesize] at (1.45,-3.95) {$\E[Y^{a=0}]$};
  % 关联：两个半菱形
  \begin{scope}[shift={(-1.45,-5.15)}]
    \fill[white] (0,0.62)--(-0.8,0)--(0,-0.62)--cycle;
    \draw[thick] (0,0.62)--(-0.8,0)--(0,-0.62)--cycle;
    \draw[thick] (0,0.62)--(0,-0.62);
  \end{scope}
  \begin{scope}[shift={(1.45,-5.15)}]
    \fill[black!28] (0,0.62)--(0.8,0)--(0,-0.62)--cycle;
    \draw[thick] (0,0.62)--(0.8,0)--(0,-0.62)--cycle;
    \draw[thick] (0,0.62)--(0,-0.62);
  \end{scope}
  \node[font=\footnotesize] at (0,-5.15) {vs.};
  \node[font=\footnotesize] at (-1.45,-6.05) {$\E[Y\mid A=1]$};
  \node[font=\footnotesize] at (1.45,-6.05) {$\E[Y\mid A=0]$};
\end{tikzpicture}
\caption{因果与关联：前者比较同一总体在“人人受治疗”与“人人未受治疗”两种情形下的结局，后者比较实际受治疗者与未受治疗者这两个子集的结局。}
\label{fig:1.1}
\end{figure}

因果的定义意味着对比整个白色菱形（所有人受治疗）与整个灰色菱形（所有人未受治疗）；而关联意味着对比原菱形中白色部分（受治疗者）与灰色部分（未受治疗者）。也就是说，关于因果的推断关心反事实世界中的“如果……会怎样”问题，例如“如果所有人都受治疗，风险会是多少？”“如果所有人都未受治疗，风险会是多少？”；而关于关联的推断关心真实世界中的问题，例如“受治疗者的风险是多少？”“未受治疗者的风险是多少？”

我们可以用迄今发展的记号把因果与关联之别形式化。风险 $\Pr[Y=1\mid A=a]$ 是一个条件概率：在满足“实际接受了治疗取值 $a$”（即 $A=a$）这一条件的总体子集中，$Y$ 的风险。相反，风险 $\Pr[Y^{a}=1]$ 是一个无条件概率——也称边缘概率：在\emph{整个}总体中 $Y^{a}$ 的风险。因此，关联由总体中两个不相交子集（由个体实际治疗取值 $A=1$ 或 $A=0$ 决定）中的风险差异来定义；而因果由同一总体在两种不同治疗取值（$a=1$ 或 $a=0$）下的风险差异来定义。在本书中，我们常用“因果效应”这一略显冗余的说法，以避免与“效应”一词常被用来仅仅表示关联相混淆。

这些截然不同的定义，解释了“关联不是因果”这句著名格言。在我们的总体中，存在关联，因为受治疗者的死亡风险（$7/13$）大于未受治疗者（$3/7$）。然而不存在因果，因为人人受治疗时的风险（$10/20$）与人人未受治疗时的风险相同。如果接受心脏移植的人平均而言比未接受者病情更重，那么因果与关联之间的这种背离就不足为奇了。在第 7 章中，我们把这种背离称为\emph{混杂}。

因果推断需要像表 1.1 那样的假设数据，但我们能指望拥有的，始终只是表 1.2 那样的真实世界数据。于是问题变成：在什么条件下，真实世界数据可用于因果推断？下一章给出一个答案：做随机化实验。

